\section{Adaptive exploration and the final contradiction}
\label{sec:exploration}

We use the scales from Lemma~\ref{prelim:scales} to construct an infinite
open cluster at the parameter $q<p_c$. Recall the coarse cubes and their
inner boxes:
\[
 c_v=(20r+1)(v_1,v_2,0),\qquad
 Q_v=\Lambda_{10r}(c_v),\qquad B_v=\Lambda_{4r}(c_v),
 \qquad v\in\mathbb Z^2.
\]
For adjacent coarse vertices $v,x$, the interface $I_{vx}$ consists of the
lattice edges between $Q_v$ and $Q_x$. Different cubes have disjoint
vertex sets. We keep the source $o=(0,0,0)$ in the root cube $Q_0$.

\subsection{The exploration and its conditional promise}

Retaining conditional connection predictions during adaptive exploration
is the method used in Kozma and Nitzan's proof of Theorem~6
\cite[Section~4, equations~(30)--(32)]{KozmaNitzan2024}.
Here each pending cube's unknown bits are reserved until that cube is processed.

Let $\mathcal I$ be the event that all bits internal to $Q_0$ are open
and, for bonds, that all bits on its four interfaces are open. This event
has positive probability because it prescribes finitely many bits and
$q>0$. Write
\[
 \mathbb P_{\mathcal I}(\,\cdot\,)
   =\mathbb P_q(\,\cdot\mid\mathcal I).
\]
Under this law, all bits not prescribed by $\mathcal I$ are still
independent Bernoulli($q$).

Declare the coarse root good and put its four neighbors in a queue, each
with parent the root. A coarse vertex is \emph{visited} once it has been
put in the queue; the root is also visited. A queued vertex is
\emph{pending} until it is processed. Use a fixed ordering of neighbors
whenever vertices are added to the queue, and always process the first
pending vertex. No vertex is assigned a parent or queued more than once.

For each pending vertex we will retain a high conditional probability
of connection from $o$. The graph used for this prediction follows
its parent chain, so that the only bits still needed to decide the
connection will be in the pending cube and its incoming interface.
Specifically, for a pending vertex $v$, let $G(v)$ consist of the
cubes on its parent chain back to the root, all edges internal to
those cubes, and the interfaces between successive cubes of the
chain. We include no other edges, even when nonsuccessive cubes
are adjacent. Set
\[
 H_v=\{o\leftrightarrow B_v\text{ in }G(v)\}.
\]
The construction will maintain the promise
\begin{equation}
 \mathbb P_{\mathcal I}(H_v\mid\text{revealed history})>1-\delta
 \quad\text{for every pending }v.
 \label{exp:promise}
\end{equation}
A cube will acquire children only after its own connection is
verified and those children satisfy the same promise.

The invariant holds initially. For a root neighbor $v$, choose a
quarter-face of $Q_v$ facing $Q_0$. Every vertex of that quarter-face has
a neighbor in $Q_0$. By \eqref{prelim:quarterface}, the probability that
$\Lambda_m(c_v)$ connects to this quarter-face inside $Q_v$ is greater
than $1-\gamma$. Here $\Lambda_m(c_v)\subseteq B_v$. Such a connection
joins $B_v$ to $o$: the root cube and the intervening edge are open in
the bond model, while both endpoints of the intervening edge are open
in the site model. Since $\gamma<\delta$, this proves
\eqref{exp:promise} for the initial queue.

To process a pending vertex $v$, reveal the bits internal to $Q_v$ and, for
bonds, its incoming parent interface, except for bits already known.
The ancestor cubes and the interfaces joining them have been
revealed at earlier steps, so this determines $H_v$.
For each unvisited neighbor $x$ of $v$, extend $G(v)$ by $Q_x$,
its internal edges, and $I_{vx}$, obtaining a graph $G'$, and compute
\begin{equation}
 b_x=\mathbb P_{\mathcal I}
 \bigl(o\leftrightarrow B_x\text{ in }G'
       \mid\text{history after this reveal}\bigr).
 \label{exp:child-probability}
\end{equation}
The number $b_x$ is calculated from the revealed values and the
product law of the remaining bits; its computation reveals no bits.
Declare $v$ good if $H_v$ holds and $b_x>1-\delta$ for every such
neighbor. In this case, give each unvisited neighbor parent $v$
and add it to the back of the queue. Otherwise declare $v$ bad and
add no vertices. Remove $v$ from the queue in either case, and
continue until the queue is empty, if this ever occurs.

\begin{lemma}[Freshness and persistence]
\label{exp:freshness}
Conditional on any positive-probability finite revealed history, all
unrevealed bits are independent Bernoulli($q$). Moreover,
\eqref{exp:promise} holds at every step of the exploration.
\end{lemma}

\begin{proof}
The initialization fixes finitely many bits and leaves the others
independent with parameter $q$. At every later step, the preceding
revealed values determine which finite set of bits is exposed.
Conditioning on the values of this selected set again leaves the
remaining bits with their product law. The classifications, queue,
parent assignments, and numbers $b_x$ are all functions of revealed
values. They therefore add no conditions on the remaining bits.
Induction proves the first assertion.

When $x$ is assigned parent $v$, its graph $G(x)$ is precisely the
graph $G'$ used to compute $b_x$. The rule for declaring $v$ good
therefore establishes \eqref{exp:promise} for $x$ at that time.
To prove persistence, it suffices to verify that processing other
vertices cannot expose any of the unknown bits of $H_x$.

All ancestor cubes and their connecting interfaces have already
been revealed. Thus the unknown bits of $H_x$ lie in $Q_x$ or,
for bonds, on $I_{vx}$. The cubes are vertex-disjoint, so another
vertex's processing does not reveal an internal bit of $Q_x$.
Each step also reveals only the processed vertex's incoming parent
interface. Distinct parent links use distinct interfaces. Moreover, a pending
vertex has no children, since only processed good vertices acquire them.
Hence $I_{vx}$ cannot be revealed by another step while $x$ is
pending. If this interface was prescribed at initialization because
$v$ is the root, its fixed values simply remain part of the history.

Consequently, until $x$ is processed, the known bits of $H_x$ stay
fixed and its unknown bits retain the same product law. Its
conditional connection probability is unchanged. Together with the
initial verification, this proves \eqref{exp:promise} at every step.
\end{proof}

\subsection{A uniform bound on bad vertices}

\begin{lemma}
\label{exp:conditional-bad}
Immediately before any pending vertex is processed, its conditional
probability of being declared bad is less than $f$.
\end{lemma}

\begin{proof}
Condition on a positive-probability history immediately before
processing $v$. Let $\mathcal F$ denote this history and
$\mathcal F^+$ the history after the prescribed reveal at $v$.
The unvisited neighbors to be tested are already determined by
$\mathcal F$.

Fix one such neighbor $x$ and the corresponding graph $G'$. Under the
conditioning on $\mathcal F$, every bit internal to the prism
$D=Q_v\cup Q_x$ has its fresh parameter-$q$ law. To check this carefully,
neither cube has been processed: $v$ is pending and $x$ is unvisited.
Their interface has not been revealed as a parent link, since that
would require $x$ to have been visited or $v$ to have been processed.
Neither vertex is the root, so their interface is not one prescribed
by initialization. Although an initialized root interface may meet
$Q_v$, it is external to $D$. Lemma~\ref{exp:freshness} gives the
required independence.

The graph $G'$ contains every edge internal to this prism, and $o$ lies
outside it in the disjoint root cube. All exterior bits have independent
laws, with probabilities $0$ or $1$ for known bits and $q$ for unknown
bits. Thus Lemma~\ref{geom:relay} applies. If
$C_x=\{o\leftrightarrow B_x\text{ in }G'\}$, it gives
\[
 \mathbb P_{\mathcal I}(H_v\cap C_x^{\mathrm c}\mid\mathcal F)
 \le 13\eta,
\]
because $H_v$ implies a connection to $B_v$ in the larger graph $G'$.
Notice that the relay estimate is applied \emph{before} revealing
$Q_v$, not after conditioning on $H_v$.

The event $H_v$ is determined by $\mathcal F^+$, and
$b_x=\mathbb P_{\mathcal I}(C_x\mid\mathcal F^+)$. The tower property
therefore yields
\[
 \mathbb E_{\mathcal I}
   [\mathbf 1_{H_v}(1-b_x)\mid\mathcal F]\le13\eta.
\]
On $H_v\cap\{b_x\le1-\delta\}$ the integrand is at least $\delta$,
so
\[
 \mathbb P_{\mathcal I}
   (H_v\cap\{b_x\le1-\delta\}\mid\mathcal F)
 \le\frac{13\eta}{\delta}.
\]
There are at most four tests. The promise and a union bound now give
\[
 \mathbb P_{\mathcal I}(v\text{ is declared bad}\mid\mathcal F)
 \le\delta+\frac{52\eta}{\delta}<f,
\]
where the last inequality is \eqref{prelim:budgets}. No independence
between the tests is needed.
\end{proof}

\begin{lemma}
\label{exp:bad-set}
For every deterministic set of $k$ distinct coarse vertices, the
probability under $\mathbb P_{\mathcal I}$ that all are processed and
declared bad is at most $f^k$.
\end{lemma}

\begin{proof}
The assertion is immediate if the set contains the root, which is
declared good. Otherwise, fix such a set $S$. For $1\le j\le k$, let
$T_j$ be the processing time of the $j$th member of $S$ to be processed,
and set $T_j=\infty$ if fewer than $j$ members are ever processed. Let
$E_j$ be the event that $T_j<\infty$ and these first $j$ members are
all declared bad; put $E_0$ equal to the whole sample space.

For each finite time $t$, the event
\[
 \{T_j=t\}\cap E_{j-1}
\]
is determined before the reveal at time $t$. Indeed, at that moment
one knows the next vertex, how many members of $S$ have already
been processed, and whether all of those members were bad.
On $\{T_j=t\}$ there are exactly $j-1$ such previous members,
so these observations determine $E_{j-1}$ there.
After termination we keep the history constant; events specifying
a later processing time are empty. Applying
Lemma~\ref{exp:conditional-bad} at time $t$ and summing gives
\begin{align*}
 \mathbb P_{\mathcal I}(E_j)
 &\le f\sum_{t\ge1}
       \mathbb P_{\mathcal I}(T_j=t,\ E_{j-1})\\
 &\le f\,\mathbb P_{\mathcal I}(E_{j-1}).
\end{align*}
Iteration yields $\mathbb P_{\mathcal I}(E_k)\le f^k$. Since no
vertex is processed twice, $E_k$ is exactly the event in the statement.
The proof does not assume that every member of $S$ is eventually
processed.
\end{proof}

\subsection{Contours and completion of the proof}

Let $X$ be the set of good coarse vertices, including the root. If $X$
is finite, only finitely many vertices can ever enter the queue: new
entries arise only from good vertices, with at most four per vertex.
The queue must therefore empty after finitely many steps. Every vertex
outside $X$ adjacent to $X$ has then been processed and declared bad.
Indeed, the root's neighbors were all queued at initialization. When
any nonroot good vertex was processed, each of its neighbors was either
already visited or assigned at that step. Thus every outside neighbor
of $X$ is processed before termination and is not good, so it is bad.

Apply Lemma~\ref{prelim:contours} to this finite set $X$, which contains
the origin of the coarse lattice. Its boundary contains a cycle of
some length $n$ in the deterministic family of at most $2n3^n$
candidates. For each candidate, fix the matching of crossed primal
edges supplied by that Lemma. It has $k\ge n/7$ edges with pairwise
disjoint endpoints. If the candidate occurs as a boundary cycle of
$X$, every one of these edges has a bad endpoint outside $X$.

For a fixed candidate there are $2^k$ ways to choose one endpoint from
each matching edge. Each choice is a deterministic set of $k$ distinct
vertices, so Lemma~\ref{exp:bad-set} bounds the probability that all
chosen vertices are processed bad by $f^k$. Hence the candidate's
probability is at most
\[
 2^k f^k=(2f)^k\le(2f)^{n/7},
\]
using $2f<1$. We have not conditioned on the random set $X$ or on the
occurrence of a cycle. Taking the union bound over the deterministic
candidate families and using the choice of $f$ in
Lemma~\ref{prelim:scales} gives
\begin{equation}
 \mathbb P_{\mathcal I}(|X|<\infty)
 \le\sum_{n\ge1}2n3^n(2f)^{n/7}<1.
 \label{exp:infinite-good}
\end{equation}

Thus, with positive probability under $\mathbb P_{\mathcal I}$,
infinitely many coarse vertices are good. Every good nonroot vertex
$v$ satisfies the actual connection event $H_v$. The boxes $B_v$ are
pairwise disjoint, so these connections put infinitely many lattice
vertices in the open cluster of $o$. Therefore
\[
 \mathbb P_q(o\text{ belongs to an infinite open cluster})
 \ge \mathbb P_q(\mathcal I)\,
       \mathbb P_{\mathcal I}(|X|=\infty)>0.
\]
This contradicts $q<p_c$. The assumption of critical percolation was
therefore false. The construction and all its estimates apply to
both bond and site percolation with their respective critical
parameters, proving Theorem~\ref{thm:main}.
